\Needspace{10\baselineskip}
\section{Finite prefixes and the arithmetic boundary}\label{sec:global-prefix}
A finite prefix is not a union of independent dyadic experiments. Its
retained bulk must be treated in one Poisson comparison, using the
macroscopic relation estimate. The starts near the origin need a different
argument: there the initial run costs one sign per prime, rather than one
sign per integer. The two candidate masses are
$\alpha_L=2^{-\pi(L)}$ at the boundary and
$\lambda_M=M2^{-L}$ in the bulk. To compare them, we must bound the other
microscopic starts, discard the intervening region, and approximate the
bulk with error $o(\lambda_M)$. A critical-scale error of size $o(1)$ is
not sufficient for this last task.

The limiting source weights arise by normalizing these two masses.
The stronger conclusion of \cref{thm:two-clock} also distinguishes their
marks: both excess laws are geometric, but one counts new primes and the
other new integer positions. This distinction is arithmetic, not a
property of the common geometric distribution.

\subsection{The exact boundary event}
Set
\begin{align*}
 C_L&=\1_{\{f(1)=\cdots=f(L)\}},&
 \mathbf B_L&=(C_L,(J_{x,L})_{2\le x\le2L^2}),\\
 B_L^{\rm mic}&=C_L+\sum_{2\le x\le2L^2}J_{x,L},&
 q_L&=\PP(B_L^{\rm mic}>0).
\end{align*}
Because $f(1)=1$, the event $C_L=1$ prescribes precisely the prime signs
up to $L$. Hence
\begin{equation}\label{eq:exact-border}
 \PP(C_L=1)=2^{-\pi(L)}
\end{equation}
exactly. The quantity $q_L$ is different: it includes all other microscopic
starts. The next result shows that their combined mass is smaller by a
fixed exponential factor in $\pi(L)$.

\begin{theorem}[Microscopic non-vacancy and localization]\label{thm:boundary}
Let $\delta_*=28271/332640$. As $L\to\infty$,
\begin{equation}\label{eq:interior-micro}
 \sum_{x=2}^{2L^2}\PP(J_{x,L}=1)
       \le2^{-(1+\delta_*-o(1))\pi(L)}.
\end{equation}
Consequently, for every fixed $0<\delta_0<\delta_*$,
\begin{equation}\label{eq:boundary-asymptotic}
 q_L=2^{-\pi(L)}(1+O_{\delta_0}(2^{-\delta_0\pi(L)})).
\end{equation}
In particular $\delta_0=1/12$ is admissible. Conditional on microscopic
non-vacancy, the exact border is the unique recorded start with probability
$1-O_{\delta_0}(2^{-\delta_0\pi(L)})$.
\end{theorem}
\begin{proof}
Let $M_x$ be the matrix whose rows are $v(n)$, $n\in V_x$.
The start matrix is $A_x=D_xM_x$, where
\[
 D_xz=(z_{x-1}+z_x,z_x+z_{x+1},\ldots,z_x+z_{x+L-1}).
\]
The kernel of $D_x$ consists of constant vectors. Therefore
\begin{equation}\label{eq:boundary-rank-drop}
 \rank A_x\ge\rank M_x-1,\qquad
 \PP(J_{x,L}=1)\le2^{-\rank A_x}.
\end{equation}
If $2\le x\le L-1$, both $2x$ and $2x-2$ lie in the proposed constant
run. The first forces $f(2)=1$, and the second forces $f(2)=-1$.
Thus these starts are impossible. For the three transition starts
$x\in\{L,L+1,L+2\}$, primes between $2\sqrt B$ and $B$, together
with prime vertices above $B$, give an identity minor of size
$(2-o(1))\pi(B)$. The explicit choice of rows is in
\cref{supp:lem:transition}.

For $B+2\le x\le B^2+2$, the uniform Laishram--Shorey bound gives
\[
 \omega\left(\prod_{n\in V_x}n\right)\ge(2-o(1))\pi(B).
\]
After discarding the prime divisors at most $B$, at least
$t_x\ge(1-o(1))\pi(B)$ primes $q>B$ remain. All vertices are at most
$B^2+B<q^2$, so such a prime occurs to valuation one at one vertex only;
two such primes cannot occupy the same vertex. Their columns span a
$t_x$-dimensional coordinate subspace.

Use the medium primes $B/11<p\le B$. Their total odd-valuation incidence
count satisfies
\begin{equation}\label{eq:medium-incidences}
 E_x\ge(H_{11}-1-o(1))\pi(B),\qquad H_{11}=\sum_{j=1}^{11}j^{-1}.
\end{equation}
Here the loss from even valuations is only $O(1)$ \emph{over the entire
prime range}: if $p^2\mid n\le B^2+B$ then $n=ap^2$ with $a\le122$,
and for each $a$ there are only $O(1)$ possibilities in an interval of
length $B$. This global count avoids paying one lost incidence per prime.

Delete the private rows. Each deleted row removes at most one medium
incidence, each remaining row has weight at most two, each column has
weight at most eleven, and a pair of medium primes cannot occur together
in two rows. The associated incidence graph is therefore simple.
The elementary simple-incidence estimate \cref{supp:lem:simple-incidence}
gives rank at least $E/12$. Passing to the quotient by the private
coordinate subspace yields
\[
 \rank M_x\ge t_x+\frac{(E_x-t_x)_+}{12}
 \ge\left(1+\frac{H_{11}-2}{12}-o(1)\right)\pi(B).
\]
Since $(H_{11}-2)/12=28271/332640$, \eqref{eq:boundary-rank-drop} proves the desired
uniform surplus in the quadratic range.

For $x>B^2+2$, the pointwise Balasubramanian--Shorey consequence
\cref{supp:lem:pointwise-maximum} gives
\begin{equation}\label{eq:gL}
 m_B(x)\le B-g_L,\qquad
 g_L=\frac B{\log B}
   (\log\log B-\log\log\log B-\theta_0+o(1)),
 \qquad \frac{g_L}{\pi(B)}\longrightarrow\infty.
\end{equation}
The affine bound then gives $\PP(J_{x,L}=1)\le2^{-g_L}$.
Sum the four ranges; their polynomial number of starts is absorbed in
the $o(\pi(B))$ exponent. Since $\pi(B)=\pi(L)+O(1)$, this proves
\eqref{eq:interior-micro}. Equation \eqref{eq:exact-border} and a union
bound prove the remaining conclusions. The complete incidence count,
the uniform literature input and the choice $K=11$ are detailed in
\cref{supp:sec:boundary}, with local product thresholds distinguished
from the thresholds of the asymptotic inputs.
\end{proof}

\subsection{Discarding deep starts without assuming independence}
\begin{lemma}[Two-defect windows at intermediate heights]
\label{lem:meso-defects}
Assume $B\asymp\log M$. The number of starts
$x\in[2L^2,2M]\cap\Z$ with at least two $B$-defective vertices is at most
$\exp(O(\log M/\log\log M))$. The same bound holds for every subpopulation;
in particular it is uniform on each $[X,2X)$ with $2L^2\le X\le M$.
\end{lemma}
\begin{proof}
Two such vertices have forms $d_1a^2,d_2b^2$ differing by a nonzero
integer of modulus at most $L$, with squarefree $B$-smooth $d_i$.
There are at most $4^{\pi(B)}=\exp(O(\log M/\log\log M))$ coefficient
pairs and $O(B^2)$ offset pairs. For $d_1\ne d_2$, use the uniform
Pell bound at ambient scale $M$; the coefficients and variables are
polynomially bounded in $M$. For equal coefficients use
$(a-b)(a+b)=(i-j)/d_1$ and the divisor bound. A solution determines the
start once the offsets are fixed. This counts the entire population up
to height $2M$ in a single Pell box, not once per dyadic interval.
\end{proof}

\begin{proposition}[Deep-start and mesoscopic estimates]\label{prop:deep-starts}
For each fixed $0<\delta<1$ and fixed logarithmic band,
\begin{equation}\label{eq:deep-starts}
 \sum_{2\le x<M^\delta}\PP(J_{x,L}=1)
 \ll \lambda_M M^{\delta-1}+\exp(-c\log M/\log\log M),
 \qquad\lambda_M=M2^{-L},
\end{equation}
where $c>0$ depends only on that band and $\delta$.
More precisely,
\begin{equation}\label{eq:mesoscopic}
 \PP(\exists x\in(2L^2,M^\delta):J_{x,L}=1)
 \le M^\delta2^{-L}+C_1 e^{C_2L/\log L}2^{-g_L}.
\end{equation}
If $L\to\infty$, $L=O(\log M)$ and $\lambda_M\to0$, the latter bound
is $o(q_L+\lambda_M)$.
\end{proposition}
\begin{proof}
For any deterministic $A\subset[2L^2,2M]\cap\Z$, starts with at most
one defect have probability at most $2^{-L}$. The remaining starts are
counted together by \cref{lem:meso-defects}, and \eqref{eq:gL} bounds each
probability by $2^{-g_L}$. Uniformly in $A$,
\[
 \sum_{x\in A}\PP(J_{x,L}=1)
 \le |A|2^{-L}+C_1 e^{C_2L/\log L}2^{-g_L}.
\]
The constants depend only on the fixed logarithmic band. Adding the
microscopic range, controlled by \cref{thm:boundary}, proves
\eqref{eq:deep-starts}; restricting this mask proves \eqref{eq:mesoscopic}.
No exterior factor counting height slices is needed. For the last assertion,
$M^\delta2^{-L}=o(\lambda_M)$, while
$g_L-\pi(L)\gg L\log\log L/\log L$ dominates the sieve prefactor.
Thus the second term is $o(2^{-\pi(L)})=o(q_L)$.
\end{proof}

The same argument gives cutoff stability. For $2L^2\le K\le2^L$,
\begin{equation}\label{eq:cutoff-stability}
 \PP(\exists x\in(2L^2,K]:J_{x,L}=1)
 \le K2^{-L}+O(e^{CL/\log L}2^{-g_L}).
\end{equation}
One may use the artificial ambient scale $2^L$ in
\cref{lem:meso-defects}. If $K=o(2^{L-\pi(L)})$, the entire bound is
$o(q_L)$. Hence the boundary asymptotic and conditional uniqueness remain
valid with cutoff $K$, and the two conditional microscopic configurations,
padded by zeros, are at total-variation distance $o(1)$ by
\cref{lem:conditioning-floor}.

\subsection{Global laws and the longest run}
\begin{theorem}[Finite-prefix Poisson approximation]\label{thm:prefix-poisson}
In a fixed logarithmic band, for every $\varepsilon>0$,
\begin{align}
 &\dTV\left(\Law\left(\sum_{2\le x<M}J_{x,L}\right),\Pois(M2^{-L})\right)
 \notag\\
 &\quad\ll\min\{1,\lambda_M(e^{-V_M+o(\nu_M)}+M^{-1/3+\varepsilon})\}
               +e^{-c\log M/\log\log M}.
               \label{eq:global-rate}
\end{align}
For $W_{M,L}=C_L+\sum_{2\le x\le M-L+1}J_{x,L}$, enlarge the
constant multiplying the existing deep-start remainder and add
\begin{equation}\label{eq:prefix-correction}
 2^{-\pi(L)}+O(L2^{-L}).
\end{equation}
In particular $\PP(R_M<L)=e^{-M2^{-L}}+o_C(1)$ at critical lengths.
\end{theorem}
\begin{proof}
Retain the single interval $U_{M,1/2}$. The proof of
\cref{thm:transfer} uses only support heights $O(M)$, the one-window
mass, the edge count and the separated-pair profile. These are supplied
by \eqref{eq:macro-defects}, \eqref{eq:edge-count} at upper scale $M$,
and \cref{thm:master-profile}. Thus \eqref{eq:hard-rate} applies on the
retained interval with mean $|U_{M,1/2}|2^{-L}$. Restoring deep starts
costs \eqref{eq:deep-starts}; the deterministic target-mean discrepancy
is $O(\lambda_M M^{-1/2})$. This proves \eqref{eq:global-rate}.

The event $W_{M,L}=0$ is exactly $R_M<L$: every long run contained in
the prefix either begins at one or has a left-maximal start at most
$M-L+1$. The full interior count differs only if the border event occurs
or a start lies in the excluded mask $M-L+2\le x<M$.
This mask has at most $L$ sites and lies above $2L^2$ for large $M$.
The mask bound in the proof of \cref{prop:deep-starts} gives total mass
\[
 L2^{-L}+O(e^{CL/\log L}2^{-g_L})
 \le L2^{-L}+O(e^{-c\log M/\log\log M}).
\]
Absorb the second term into the remainder already present in
\eqref{eq:global-rate}. The exact border probability then proves
\eqref{eq:prefix-correction} and the void law.
\end{proof}

\begin{corollary}[Almost-sure asymmetric envelopes]\label{cor:as-longest}
For every fixed $\varepsilon>0$, with probability one, eventually
\begin{align}
 R_M-\log_2M&\ge-\log_2(\log\log M)-O(1),
                                          \label{eq:as-lower}\\
 R_M-\log_2M&\le\log_2(\log M)
        +(1+\varepsilon)\log_2(\log\log M)+O_\varepsilon(1).
                                          \label{eq:as-upper}
\end{align}
In particular $R_M=\log_2M+O(\log\log M)$ almost surely.
\end{corollary}
\begin{proof}
Set $M_k=2^k$ and, for sufficiently large $k$, choose
\[
 L_k^-=\lfloor k-\log_2(\log k)-1\rfloor,\qquad
 L_k^+=\lceil k+\log_2k+(1+\varepsilon)\log_2(\log k)\rceil.
\]
The lower intensity satisfies $2\log k\le M_k2^{-L_k^-}<4\log k$,
so the target void probability is at most $k^{-2}$. The upper intensity
is at most $1/[k(\log k)^{1+\varepsilon}]$, a summable sequence.
Every error in \eqref{eq:global-rate}--\eqref{eq:prefix-correction} is also
summable: the cutoff error is a polynomial in $k$ times
$e^{-c\sqrt{k\log k}}$, and the deep-start and border errors are at most
$e^{-c'k/\log k}$. Borel--Cantelli gives
$L_k^-\le R_{2^k}<L_k^+$ for all sufficiently large $k$, almost surely.
For $2^k\le M<2^{k+1}$ use $R_{2^k}\le R_M\le R_{2^{k+1}}$.
The change from $k$ to $\log M/\log2$ and from $k$ to $k+1$ costs only
a bounded additive term. This proves both envelopes. No independence
between different scales has been used, and no matching lower bound for
the upper fluctuations is asserted.
\end{proof}

\begin{theorem}[Macroscopic exact-mark transport]\label{thm:macro-field}
Let $a_M=o(\log M)$ be a nonnegative integer, $b_M=\lfloor\log_2M\rfloor$,
$L_M=b_M-a_M$, and $\lambda=M2^{-L_M}$. Define
\[
 \mathcal E^{\rm base}_{M,a_M}
 =\sum_{x=2}^{M-L_M+1}\sum_{r\ge-a_M}
             K^{\rm ex}_{x,b_M+r}\delta_{(x/M,r)}.
\]
Its target has independent Poisson masses $2^{-b_M-r-1}$ on the same
lattice, and total intensity $(M-L_M)2^{-L_M}$.
For $A_M\in\mathcal F_{Y_M^\sharp}$ with $\PP(A_M)>0$, put
$I_M^{\rm info}=-\log\PP(A_M)$;
the full-field distance is bounded by
\begin{equation}\label{eq:macro-field-rate}
 C_\varepsilon e^{I_M^{\rm info}}
 \left[\lambda(1+\lambda)(e^{-V_M+o(\nu_M)}+M^{-1/3+\varepsilon})
            +\lambda M^{-1/2}+e^{-c\log M/\log\log M}\right].
\end{equation}
It tends to zero under the hard labelled budget
$I_M^{\rm info}+\log\lambda+\log(1+\lambda)\le V_M-c'\nu_M$.
The movable-cutoff labelled budget and the aggregated one-factor budget
of \cref{thm:moving-comparison} also hold, with $N$ replaced by $M$ and
exact-level target means $(M-L_M)2^{-b_M-r-1}$.
\end{theorem}
\begin{proof}
On $[\lceil M^{1/2}\rceil,M-L_M+1]$ use the finite marked comparison
with the macroscopic master profile; the identities
\eqref{eq:mark-folding} still collapse the separated-pair term to the
base threshold event. Restore lower starts with \eqref{eq:deep-starts}
and condition using \cref{lem:stable-lift}. Remove the exact-mark tail
using \eqref{eq:mark-tail} and the global scalar bound. A cutoff
$E=\lceil(I_M^{\rm info}+V_M+\log(2+\lambda))/\log2\rceil$ is
$o(\log M)$ whenever the bound is nontrivial; the deep-start tail error
is absorbed by the last term of \eqref{eq:macro-field-rate}.
This proves the displayed bound.
For the movable cutoff repeat with $w=\widehat V_M$.
For aggregation, sum over positions before applying the directional Stein
factor; the deep-start error is negligible under its information
budget. \Cref{supp:sec:macro-marked} details the finite-cutoff calculation.
\end{proof}
The superscript ``base'' is important: only the base $L_M$-window is
required to be contained. The exact right endpoint of a larger marked run
may lie beyond $M$. This is not the list of right-censored maximal runs of
the finite sequence. All measurable functionals of this base-contained
field, including spatial masks and order statistics, transfer under the
labelled comparison. Aggregated couplings do not retain locations.

\begin{figure}[htbp]\centering
\begin{tikzpicture}[x=.85cm,y=1cm]
\draw[->] (0,0)--(13,0);
\draw[thick] (1,0)--(11,0);
\foreach \x/\t in {1/1,6/x,11/M}{\draw (\x,.09)--(\x,-.09);
 \node[below,font=\small] at (\x,-.1) {$\t$};}
\draw[draw=FigBlue,line width=2pt] (6,.45)--(9.5,.45);
\node[above,font=\small,text=FigBlue] at (7.75,.45) {contained base window};
\draw[draw=FigRust,dashed,line width=1.2pt] (6,1.25)--(12,1.25);
\node[above,font=\small,text=FigRust] at (9,1.25) {exact run may continue beyond $M$};
\draw[dotted] (11,-.1)--(11,1.6);
\end{tikzpicture}
\caption{Containment of a base window does not censor an exact mark.
The longest-run void event instead uses the threshold-specific upper
start $M-L+1$. The two conventions agree where stated, but are not
interchangeable.}\label{fig:containment}
\end{figure}

\subsection{Relative bulk approximation and the crossover}
Assume from now on $L\to\infty$, $L=O(\log M)$ and
$\lambda_M=M2^{-L}\to0$. Fix $0<\delta<1$ and put
$H_M=\lceil M^\delta\rceil$,
$I^{\rm bulk}_{M,L}=[H_M,M-L+1]\cap\Z$ and
$\Xi^{\rm bulk}_{M,L}=\sum_{x\in I^{\rm bulk}_{M,L}}J_{x,L}\delta_{x/M}$.
Let $\Pi^{\rm bulk,lat}_{M,L}$ have independent Poisson masses $p=2^{-L}$
at those same sites; its intensity is
$\lambda_M^{\rm bulk}=|I^{\rm bulk}_{M,L}|p\sim\lambda_M$.
Fix an upper-band constant $A_0$ with $L\le A_0\log M$ for all large
$M$. The rare-intensity assumption supplies the positive lower band.
Constants and little-oh terms in this subsection may depend on
$\delta$, $A_0$ and the displayed error exponent, but not on $M$ or $L$.

\begin{theorem}[Relative signed marked bulk kernel]\label{thm:relative-bulk}
Let $\mathcal K^{\rm bulk}_{M,L}=(K^s_{x,e})_{x\in I^{\rm bulk}_{M,L},\,e\ge0,\,s\in\{\pm1\}}$
and let $\mathcal P^{\rm bulk}_{M,L}$ have independent Poisson coordinates
of means $2^{-L-e-2}$ on this same index set. At $Y=Y_M^\sharp$, for every
fixed $0<\varepsilon<1/3$,
\begin{equation}\label{eq:relative-bulk-rate}
 \EE\,\dTV\left(\Law(\mathcal K^{\rm bulk}_{M,L}\mid\mathcal F_Y),
                         \Law(\mathcal P^{\rm bulk}_{M,L})\right)
 \ll_{\delta,A_0,\varepsilon}\lambda_M
               \left(e^{-V_M+o(\nu_M)}+M^{-1/3+\varepsilon}\right).
\end{equation}
In particular, summing marks and signs gives
\begin{equation}\label{eq:relative-bulk}
 \EE\,\dTV(\Law(\Xi^{\rm bulk}_{M,L}\mid\mathcal F_Y),
                    \Law(\Pi^{\rm bulk,lat}_{M,L}))=o(\lambda_M).
\end{equation}
\end{theorem}
\begin{proof}
For a finite mark cutoff $E$, repeat the signed comparison of
\cref{thm:marked-field} on this single macroscopic interval. With
$Q=L+E+1$ and $p=2^{-L}$ its error is at most
\[
 Cp(\mathcal M_B+|D_Y(Q)|)
       +Cp^2(MQ+E_Y(Q)+R_{2,\delta}(M,L)).
\]
The signed marks are summed before the separated-pair bound, so the
relation sum is at the base length $L$. Divide by $Mp=\lambda_M$.
The defect and support terms become $M^{-1/2+o(1)}$ and
$e^{-V_M+o(\nu_M)}$, respectively. The edge term becomes
$\lambda_M(e^{-V_M+o(\nu_M)}+M^{-1+o(1)})$.
By \eqref{eq:master-envelope} and $2^B=2/p$, the relation term is
$O_\varepsilon(M^{-1/3+\varepsilon}(1+\lambda_M))$.
The local term $pQ$ is smaller because $\lambda_M\to0$.

Take $E=\lceil V_M/\log2\rceil$. Then $E=o(\log M)$ and the base and
maximal support lengths stay in a fixed logarithmic band. The target
tail has mass at most $\lambda_M2^{-E-1}$. The actual tail is the event
that some retained site starts a run at length $L+E+1$; the macroscopic
masked first-moment bound gives the same upper bound times
$1+M^{-1/2+o(1)}$. Couple both tails to zero and average. Since
$\lambda_M\le1$ for large $M$, the preceding bounds give
\eqref{eq:relative-bulk-rate}. Measurable contraction proves
\eqref{eq:relative-bulk}.
\end{proof}
Every coordinate of $\mathbf B_L$ depends only on prime signs up to
$2L^2+L<Y_M^\sharp$. The stable product lift therefore gives
\begin{equation}\label{eq:boundary-bulk-factorization}
 \dTV\left(\Law(\mathbf B_L,\Xi^{\rm bulk}_{M,L}),
       \Law(\mathbf B_L)\otimes\Law(\Pi^{\rm bulk,lat}_{M,L})\right)
       =o(\lambda_M)=o(q_L+\lambda_M).
\end{equation}
This is asymptotic factorization at the indicated error scale, not exact
finite-$M$ independence.

\begin{theorem}[Two-source rare tail and conditional location]
\label{thm:crossover}
Under the preceding rare-intensity assumptions,
\begin{equation}\label{eq:two-source-tail}
 \PP(R_M\ge L)=q_L+\lambda_M+o(q_L+\lambda_M).
\end{equation}
Let $X^*_{M,L}$ be the least start of a contained run of length at least
$L$, on the event that one exists. If
$L-\log_2M-\pi(L)\to s\in\mathbb R$, then
\begin{equation}\label{eq:crossover-location}
 \Law(X^*_{M,L}/M\mid R_M\ge L)
 \Longrightarrow \frac1{1+2^{-s}}\delta_0+
                         \frac{2^{-s}}{1+2^{-s}}\Leb_{[0,1]}.
\end{equation}
The source-resolved version assigns a microscopic start its location divided
by $L^2$, and a bulk start its location divided by $M$; its limit is the
same mixture with the two source labels retained. The weights also have
their evident degenerate limits when $s\to\pm\infty$.
\end{theorem}
\begin{proof}
The mesoscopic region has mass $o(q_L+\lambda_M)$ by
\cref{prop:deep-starts}. Under the product target in
\eqref{eq:boundary-bulk-factorization}, a simultaneous microscopic and
bulk hit has probability $O(q_L\lambda_M)$, and at least two bulk hits
have probability $O(\lambda_M^2)$. Both are negligible on the combined
scale. This proves \eqref{eq:two-source-tail}.
Given a microscopic hit, \cref{thm:boundary} puts the start at one with
probability $1-o(1)$. Given exactly one bulk target point, its location
is uniform on the common lattice. Finally
\[
 \frac{\lambda_M}{q_L}
   =2^{-(L-\log_2M-\pi(L))}(1+o(1))\longrightarrow2^{-s}.
\]
Divide the two one-hit masses by \eqref{eq:two-source-tail} and pass from
the uniform grid to Lebesgue measure weakly. All total-variation comparisons
precede this final diffuse limit.
\end{proof}

\subsection{The prime clock}\label{sec:prime-clock}
Let $p_1=2<p_2<\cdots$. Almost surely
$J_\partial=\min\{j:f(p_j)=-1\}$ is finite and the maximal initial
positive run has length $\ell_\partial=p_{J_\partial}-1$: every smaller
integer has only positive prime factors, whereas that negative prime
changes sign. On this full-measure event and $C_L=1$, put
$G_L=J_\partial-\pi(L)-1$; assign $G_L=0$ elsewhere whenever a globally
defined random variable is needed. Then
\begin{equation}\label{eq:prime-clock}
 \PP(G_L=k\mid C_L=1)=2^{-(k+1)},\qquad
 \PP(\ell_\partial-L\ge h\mid C_L=1)
                    =2^{-[\pi(L+h)-\pi(L)]}.
\end{equation}
Thus prime rank, not integer distance, is its geometric clock.
The boundary theorem makes the law of $G_L$ given $B_L^{\rm mic}>0$
$O(2^{-\delta_0\pi(L)})$-close to $\nu_{\rm geo}$.
The prime number theorem gives $\ell_\partial/L\to1$ conditionally, but
the unscaled integer overshoot is not tight: for $L=m!+1$ the following
$m-1$ integers contain no prime.

\subsection{A lattice crossover with two geometric clocks}\label{sec:two-clock}
The two branches have the same geometric law but different clocks:
additional integers in a bulk run, and additional positive prime signs at
the initial boundary. A moving discrete target states their joint law
without requiring a limiting ratio of the two rare masses.

Write $\alpha_L=2^{-\pi(L)}$ and $\lambda^{\rm bulk}=\lambda_M^{\rm bulk}$.
On $\{R_M\ge L\}$, let $X^*=X^*_{M,L}$ and define
\[
 \Gamma_{M,L}=
 \begin{cases}
  (\partial,1,+1,G_L),&X^*=1,\\
  (\mathrm{bulk},X^*,f(X^*),\ell_{X^*}-L),&X^*\in I^{\rm bulk}_{M,L},\\
  \dagger,&\text{otherwise}.
 \end{cases}
\]
Here $\dagger$ records an exceptional start in the intervening region.
Let $\nu_\pm=(\delta_{-1}+\delta_{+1})/2$, and let
$\operatorname{Unif}(I^{\rm bulk}_{M,L})$ be the uniform law on that
integer grid.

\begin{theorem}[Sign-resolved, two-clock crossover]\label{thm:two-clock}
Under the rare-intensity assumptions of \cref{thm:crossover},
\begin{equation}\label{eq:two-clock-tv}
 \dTV\left(\Law(\Gamma_{M,L}\mid R_M\ge L),Q_{M,L}\right)\longrightarrow0,
\end{equation}
where the deterministic probability law on the same discrete state space is
\begin{align}
 Q_{M,L}
 &=\frac{\alpha_L}{\alpha_L+\lambda^{\rm bulk}}
        \delta_{(\partial,1,+1)}\otimes\nu_{\rm geo}
       \notag\\
 &\quad+\frac{\lambda^{\rm bulk}}{\alpha_L+\lambda^{\rm bulk}}
        \delta_{\mathrm{bulk}}\otimes
        \operatorname{Unif}(I^{\rm bulk}_{M,L})\otimes\nu_\pm
                     \otimes\nu_{\rm geo}.
       \label{eq:two-clock-target}
\end{align}
No convergence of $\lambda^{\rm bulk}/\alpha_L$ is needed. In particular,
\begin{equation}\label{eq:crossover-positive-sign}
 \PP(f(X^*)=+1\mid R_M\ge L)
  =\frac{\alpha_L+\lambda^{\rm bulk}/2}
         {\alpha_L+\lambda^{\rm bulk}}+o(1).
\end{equation}
The bulk excess in \eqref{eq:two-clock-tv} may instead be measured using
the right-censored run length in the finite word $(f(1),\ldots,f(M))$.
\end{theorem}
\begin{proof}
Fix $K\ge1$. Count the consecutive positive signs among the first $K$
primes after $L$, capped at $K$, and set the result to zero off $C_L=1$.
This is a genuine finite-cylinder random variable, agreeing almost surely
with $G_L\wedge K$. For fixed $K$, Bertrand's postulate places these
primes below $2^K L<Y_M^\sharp$ for large $M$. The stable lift applied to
\cref{thm:relative-bulk} therefore retains the microscopic vector and
this clock jointly with \emph{all} signed bulk marks, at error
$o(\lambda_M)$. No bulk-mark truncation is needed here.
The interior microscopic and mesoscopic errors are
$o(\alpha_L+\lambda_M)$ by \cref{thm:boundary,prop:deep-starts}.

Write $\alpha=\alpha_L$ and $b=\lambda^{\rm bulk}$.
In the product target, keep either a boundary hit with no bulk point,
or no boundary hit with exactly one bulk point. This event has exact mass
\[
 e^{-b}\bigl[\alpha+(1-\alpha)b\bigr].
\]
Its conditional source weights are
$\alpha/[\alpha+(1-\alpha)b]$ and
$(1-\alpha)b/[\alpha+(1-\alpha)b]$.
Their mixture differs from the weights in \eqref{eq:two-clock-target}
by at most
\[
 \frac{\alpha^2b}{(\alpha+b)[\alpha+(1-\alpha)b]}\le b.
\]
The discarded simultaneous and multiple hits cost
$O(\alpha b+b^2)=o(\alpha+b)$.
Given the one bulk point, its site, sign and excess have the independent
laws in \eqref{eq:two-clock-target}. Normalize all comparison errors by
$\alpha+b$, not by the smaller source mass. This proves the comparison
with the boundary clock capped at $K$, uniformly in the moving weights.

Under the actual law, a recorded boundary tail $G_L>K$ has conditional
probability at most $2^{-K-1}$, since
$\PP(C_L=1,G_L>K)=\alpha2^{-K-1}$ and
$\PP(R_M\ge L)\ge\alpha$. The target boundary tail is also at most
$2^{-K-1}$. Let $M\to\infty$ first and then $K\to\infty$.
This proves \eqref{eq:two-clock-tv}; projecting the sign gives
\eqref{eq:crossover-positive-sign}.

For the last assertion, under the one-hit bulk target a mark is censored
at site $x$ only if $e>M-x+1-L$. The average of this geometric tail is
\[
 \frac1{|I^{\rm bulk}_{M,L}|}
 \sum_{x\in I^{\rm bulk}_{M,L}}2^{-(M-x+2-L)}
       \le\frac1{|I^{\rm bulk}_{M,L}|}=o(1).
\]
The total-variation comparison transfers this estimate, so replacing that
mark changes the conditional law by $o(1)$.
\end{proof}
After scaling sites by $M$, convergence of the ratio recovers the mixture
in \cref{thm:crossover}, now with signs and both overshoot clocks retained.
The spatial limit is weak; the stronger statement
\eqref{eq:two-clock-tv} is entirely on the prelimit lattice. Its $o(1)$
error is on the combined source scale. It does not supply a relative
estimate for the smaller source when its mixture weight tends to zero.

\subsection{Affine information at the boundary}
\begin{theorem}[Affine-conditioned crossover]\label{thm:affine-crossover}
Let $A_M=\{G_M X(Y_M^\sharp)=b_M^{\rm aff}\}$ be a compatible affine
cylinder of rank $r_M$, and assume compatibility also with $C_L=1$.
Set
\[
 U_L=\Span\{e_p:p\le L\},\qquad
 \kappa_M=\dim(\operatorname{row}G_M\cap U_L),\qquad
 d_M=\pi(L)-\kappa_M.
\]
Whenever $L\le Y_M^\sharp$, the finite identity
$\PP(C_L=1\mid A_M)=2^{-d_M}$ holds exactly.
For the asymptotic conclusions, let $M,L\to\infty$ with
$L=O(\log M)$ and $\lambda_M=M2^{-L}\to0$.
Fix $c>0$ and assume $r_M\log2\le V_M-c\nu_M$ eventually.
The affine equations, their number and their right-hand sides may all
vary along the sequence. Then the rare-tail formula is
\[
 \PP(R_M\ge L\mid A_M)=2^{-d_M}+\lambda_M+o(2^{-d_M}+\lambda_M).
\]
If $L-\log_2M-d_M\to s$, the conditional location law given
$A_M$ and $R_M\ge L$ is \eqref{eq:crossover-location} with the same
weights. Microscopic non-vacancy under $A_M$ is concentrated on the border.
\end{theorem}
\begin{proof}
The stacked cylinder and border system has rank
$r_M+\pi(L)-\kappa_M$, proving the exact mass ratio.
Write $I=r_M\log2$. Apply the factor $e^I$ to the quantitative bound
\eqref{eq:relative-bulk-rate} \emph{before} taking a limit. Its normalized
cutoff term is at most $e^{-c\nu_M+o(\nu_M)}$, and its polynomial term
vanishes because $I=O(V_M)=o(\log M)$. Hence the conditional bulk error
is $o(\lambda_M)$.
The microscopic interior mass, after conditioning and division by
$2^{-d_M}$, is at most
$O(e^I2^{-\delta_0\pi(L)-\kappa_M})=o(1)$, since
$I=O(V_M)=o(\pi(L))$. Similarly the two terms of
\eqref{eq:mesoscopic} become $o(\lambda_M)$ and $o(2^{-d_M})$.
The conditional stable lift and the exact one-hit normalization in the
proof of \cref{thm:two-clock} now apply, with total normalization
$2^{-d_M}+\lambda_M^{\rm bulk}$.
\end{proof}
The parameter $\kappa_M$ counts linear consequences supported wholly on
the border coordinates, not input equations merely mentioning them.
No condition on future primes is needed for the rare tail, microscopic
concentration or location laws. The geometric future mark requires a
separate condition. Let $q_1<q_2<\cdots$ be the primes after $L$.
If, for each fixed $K$, eventually in $M$ with a threshold that may depend
on $K$,
\[
 (\operatorname{row}G_M+U_L)\cap\Span\{e_{q_1},\ldots,e_{q_K}\}=\{0\},
\]
then the geometric boundary mark and the full mixture of
\cref{thm:two-clock} survive under $A_M$, with $\alpha_L$ replaced by
$2^{-d_M}$. First fix $K$ and pass to the size limit, then let $K\to\infty$;
no common threshold for all $K$ or neutrality for a growing $K(M)$ is
assumed. In general, if the extended affine stack is compatible,
$\PP(G_L\ge K\mid A_M,C_L=1)$ equals
$2^{-[K-\dim((\operatorname{row}G_M+U_L)\cap\Span\{e_{q_1},\ldots,e_{q_K}\})]}$;
if incompatible, the tail is zero. This separates the information that
changes the border mass from information that changes its future mark.
\Cref{supp:sec:boundary-conditioning} records the rank calculation.
